% legends.tex : Figure legends (Fig. 1 to 7) and the Table 1 caption.
% Body text only, no preamble. Sci Rep places legends after the references.
% Spec: paper/manuscript/MANUSCRIPT_SPEC.md sections 7 and 11; panel content and
% first sentences from the v3 figure spec (v3_restructure/FIGURE_SPEC_v3.md)
% (legend drafts in sections 3.6 to 10.4). First sentences are kept as in the spec.
% Limits: each figure legend <= 350 words after the X placeholder is replaced
% (spec budget: Fig 1, 4, 6, 7 reserve 40 to 50 words for the placeholder);
% Table 1 caption <= 150 words. Count on the rendered PDF (pdftotext, wc -w).
% Numbers: read from the source tables named in paper/claims/*/README.md and
% rounded with style guide 4.5 from the unrounded file values (|delta| < 1 cm two
% decimals, other cm values one decimal, correlations two decimals, percentages
% integer, CI ends with the digits of the point estimate). Replace with
% numbers.tex macros once the number-building script exists (QA item M-09).
% Sign and symbols: U+2212 for negative numbers and U+00B0 for degrees, as in the
% sibling section files; numeric ranges use "--"; ranges with a negative end use
% "to". Panel references use (a), (b); letter ranges use "to".
% Placeholders [X?: ...] are the registered strings of spec section 11, copied
% character for character. They contain Korean text, so a build that keeps them
% needs xeCJK. [미확인: ...] marks facts not yet confirmed (spec, figure spec 14).
% Citation keys: paper/manuscript/en/refs_legends.bib (keys follow references/INDEX.md).

\section*{Figure legends}

% ---- Figure 1 -------------------------------------------------------------
% Sources: D README 2 (licence counts, source coefficient, ERA5-Land 2015-2020
% period, caution 5); figure spec 3.3 (panels, 400-cell sample, white circles, map
% elements); final batch plan 2.8 (validation ladder: three regions, Stefan fitted
% on the training fold). Central meridian of panel a: fig1.py uses the Lena Delta zoom centre
% (LON0_A = None -> lon_c 126.698 E, Fig1_values.txt), not 45 W (checked 2026-10-04).
\noindent Figure 1. Study regions, Stefan coefficients and evaluation design.
(a) Labelled cells aggregated to 0.5° blocks; circle area is proportional to the number of 1~km label locations per block, and white circles mark regions added after the label tables were fixed (Central Russia, expanded edition; Tibetan Plateau, inset; Table~1).
Shading marks continuous (permafrost fraction $\geq$90\%) and discontinuous (50--90\%) permafrost [미확인: 영구동토 구역 자료 이름과 판]; the rectangle marks panels c and d.
(b) Stefan coefficient $E = \mathrm{ALT}/\sqrt{\mathrm{TDD}}$ of each labelled cell on a logarithmic axis, where ALT is the active-layer thickness and TDD the thawing degree-days of 2~m air temperature (ERA5-Land\cite{munozsabater2021_era5_land}, 2015--2020 mean); at most 400 cells are shown per region, and Source coefficient marks the coefficient $E_0$ fitted on the source pool when the region is held out.
(c) Region holdout: the Lena Delta cells and all cells within 100~km of them are removed from the training data.
(d) Label blocks and scoring blocks of one of five splits, with 40 label cells drawn from the label blocks; errors are scored only in scoring blocks.
Within-region tests use such splits without source data and were designed after the transfer results had been viewed.
(e) Root-mean-square error (RMSE) of direct machine learning (ML) and of the Stefan model fitted to the training cells (Recalibrated Stefan) under random-cell, site, 0.5° block, $k$-fold nearest-neighbour distance matching (kNNDM)\cite{linnenbrink2024_knndm} and region-holdout validation, against the median distance from scoring cells to the nearest training cell (mean of Alaska, the Lena Delta and Canada, with regional values); [XH: Fig 1e | 결과 열람 뒤 설계, 서술(판정어 없음) | 2.8 사전 고정 서술 규칙].
Cells without a verified licence (Canada 38, North Atlantic 19, W Russia 1) are not drawn.
Polar stereographic projection (central meridian 127°~E), scale bar true at 70°~N; Tibetan Plateau inset in a Lambert azimuthal equal-area projection; coastlines, Natural Earth 50~m; Cartopy [미확인: 판].

\medskip
% ---- Figure 2 -------------------------------------------------------------
% Sources: C2 README 2.1 (four-region recalibration and residual rows at ten
% labels; uncorrected P 0.034 from the abstract bundle, column p_two), 2.2 (Canada and W Russia rows); C3 README 2.5; figure table
% fig2_region_curves (W Russia residual curve at ten labels, -12.52); resamples
% from the pool curve meta (10,000) and the regional curve table (1000). Both CIs
% agree for the four-region recalibration row and the Canada and W Russia rows.
\noindent Figure 2. With ten labels, Stefan recalibration lowered the four-region mean error; no further difference was established for residual ML.
Error change is the root-mean-square error (RMSE) of a method minus that of the source-coefficient Stefan model (zero line) on the scoring blocks of each held-out region; negative values mean lower error.
Recalibrated Stefan refits the coefficient $E$ with the $n$ target labels, shrunk towards the source coefficient ($\kappa = 10$); Anchor + residual ML adds a CatBoost residual with weight 0.25; Direct ML is CatBoost without a physics anchor; Year-matched Stefan uses thawing degree-days of the label years.
(a) Mean of the four main regions (the Lena Delta, Canada, W Russia, E Russia), up to ten labels.
(b) Mean of the Lena Delta and Canada, the regions with at least 40 candidate labels; All, every label; Year-matched Stefan is a point estimate.
(c to f) Single regions; label counts above the candidate pool (14--17 cells in W and E Russia, 325--406 in Canada) are masked, and panel e has its own vertical range (values to −12.5~cm).
Lines are cell-weighted means over five splits; bands are 95\% block-bootstrap confidence intervals (CIs) of the two recalibration-based methods (10,000 resamples for the fixed-composition means in panels a and b; 1000 in panels c to f).
With ten labels, recalibration lowered the four-region RMSE by 2.5~cm (95\% CI 1.9 to 3.0~cm; Holm-adjusted $P = 0.001$), with lower error in W Russia only (12.0~cm; 95\% CI 10.5 to 13.6~cm) and higher error in Canada (2.3~cm; 95\% CI 0.7 to 2.9~cm); adding the residual changed it by −0.18~cm (95\% CI −0.53 to −0.01~cm; $P = 0.034$ before correction, Holm-adjusted $P = 0.136$).
Verdicts of all plotted contrasts are in Supplementary Table~S2.
The label-grid analyses were internally pre-registered and run on one computing platform.

\medskip
% ---- Figure 3 -------------------------------------------------------------
% Sources: C1 README 2.3 (shuffled and linear placebo rows; abstract-bundle Holm P);
% C3 README 2.1 and 2.2 (structure rows at 0, 3, 10, 40, 160 labels; the 40-label
% physics-output row depends on the split-independence assumption), 2.6 (least
% squares minus shrinkage, descriptive). Plotted series: figure table
% fig3_b_L10_alln (physics outputs, recalibrated inputs). Resamples 10,000.
% Panel a: the rendered figure (fig3.py CONCEPT_REGION, Fig3_values.txt [6]-[7]) uses W Russia (31 cells, E0 1.589510,
% ten-label recalibrated E 2.109580), not the Lena Delta of the spec draft; the legend follows the rendered figure.
\noindent Figure 3. Physics pseudo-labels lowered error relative to shuffled pseudo-labels, and the better combination structure depended on the number of labels.
(a) Concept on the 31 labelled cells of W Russia: active-layer thickness against the square root of thawing degree-days, with the source-coefficient Stefan line (Source Stefan, 1.59~cm~(°C~d)$^{-1/2}$), the line recalibrated with ten drawn labels (Recalibrated Stefan, 2.11) and the departures of those labels that residual ML learns (Residual); thickness increases downwards.
(b) Direct ML trained with physics pseudo-labels minus the same model trained with placebo pseudo-labels, four-region mean, without labels and with ten labels.
(c) Residual structure minus physics-input structure against the number of labels, for two input sets (Physics inputs, outputs of two physics models; Recalibrated inputs, source and recalibrated Stefan predictions); the key separates registered label counts (0 and 10) from auxiliary ones, and values at 40 and 160 labels are Lena Delta and Canada means.
(d) Local least-squares minus shrinkage recalibration of the Stefan coefficient (descriptive; positive values favour shrinkage), four-region mean up to ten labels, Lena Delta and Canada mean beyond.
Intervals are 95\% block-bootstrap confidence intervals (CIs; 10,000 resamples of 0.5° scoring blocks), weighted by cell and by block as indicated in the key (panel d, by cell); the grey band marks ±0.5~cm.
In panel b, physics pseudo-labels lowered error relative to shuffled pseudo-labels by 1.6~cm without labels (95\% CI 1.0 to 2.0~cm; Holm-adjusted $P = 0.001$); the difference from linear thaw-index pseudo-labels was 0.48~cm without labels and within ±0.5~cm with ten labels.
In panel c, the residual structure had 2.5--3.7~cm lower error under both weightings with up to ten labels (Holm-adjusted $P = 0.001$ for physics inputs at ten labels); at 40 and 160 labels the physics-input structure had 1.9--2.4~cm lower error, a difference that came from Canada (one of four contrasts depends on the split-independence assumption).
The placebo and structure analyses were internally pre-registered.

\medskip
% ---- Figure 4 -------------------------------------------------------------
% Sources: C2 README 2.1 (rank correlation 0.38 [0.17, 0.55], 28 targets, signed
% -0.01), 2.2 Tibetan rows (bias 227.71, gain 159.55); C5 README 2.1 and 2.2
% (selection minus fixed recipe at 10, 40, 160 and all labels; regional rows);
% 1.2 N4 (40-label reduction and all-label non-inferiority depend on the
% split-independence assumption). Placeholder: spec section 11.
\noindent Figure 4. The error reduction of label-based correction was rank-correlated with the source-coefficient bias measured with ten labels.
Intervals in panels b and c are 95\% block-bootstrap confidence intervals (CIs; 10,000 resamples of 0.5° scoring blocks), weighted by cell and by block as indicated in the key; the grey band marks ±0.5~cm.
(a) Smallest tested label count $n^{*}$ (line from the previous grid value; Methods) at which recalibration lowered error relative to the source-coefficient Stefan model (Recalibration) and residual ML relative to the recalibrated model (Residual ML); Not reached marks targets without $n^{*}$.
Rows follow the ten-label bias; sub-regions are not independent.
(b) Error change with all labels, relative to the source-coefficient Stefan model, of the method chosen by cross-validation within the target labels, against the absolute mean bias of that model at ten labels (Spearman $\rho = 0.38$; 95\% CI 0.17 to 0.55; 28 targets; signed bias, $\rho = -0.01$); the Tibetan Plateau lies outside the axes (bias 227.7~cm, error change −159.6~cm).
(c) Cross-validation selection minus a fixed residual recipe (residual weight 0.25), four main regions at 10 and all labels, the Lena Delta and Canada at 40 and 160 labels.
The selection was non-inferior (margin 0.5~cm) at 40 and 160 labels and had 1.3~cm lower error at 160 labels (95\% CI 0.7 to 1.7~cm); the 0.87~cm reduction at 40 labels and the all-label non-inferiority depend on the split-independence assumption, and non-inferiority did not hold at ten labels.
The change came from Canada (−2.1 and −2.7~cm; Lena Delta +0.35 and +0.07~cm).
(d) Error change from recalibration and from residual ML beyond recalibration with all labels, against the same bias; [XA: Fig 4d | 사후 분석(재현(비맹검)) | 2.1 결과 범주(양, 확인하지 못함, 음)와 CE1, CE2 판].
Panels b and c were designed after earlier results had been viewed; panel d is post hoc.

\medskip
% ---- Figure 5 -------------------------------------------------------------
% Sources: C6 README 1.3, 2.1, 2.2 (within-region strategy rows), 2.5 (transfer
% rows). Common-resampling verdicts read from the within-region and transfer
% placement rows of the source tables named in C6 README 2.0: Lena block
% stratification at 100 labels, Lena active selection at 200 labels and Canada
% covariate spread at 40 transfer labels depend on the split-independence
% assumption (not recorded in C6 README). The Alaskan sub-region rows of the spec
% draft are not plotted in panel d and are left to the Supplementary Information.
% Rendered Fig5 (fig5.py, Fig5_values.txt): maps a-c (Random, Block stratified, Covariate spread; no active-selection map,
% spec decision D-12 not approved), d = budget curves, e = transfer forest. Panel letters here follow the rendered figure.
\noindent Figure 5. Spreading labels across blocks and covariate space lowered error versus random cells in Canada, not in the Lena Delta.
(a to c) One draw of 100 labels in Canada (within-region test, split 1) by the three placement procedures named above the maps (Methods); small dots are candidate cells in the label blocks, and circle area is proportional to the labels selected per site (blocks closer than 40~km pooled).
The maps illustrate the procedures and do not rank sites.
(d) Error change relative to random cells for Anchor + residual ML (residual weight 0.25) against the number of labels within Canada, the Lena Delta and Alaska.
(e) The same contrast in transfer with ten and forty labels.
Intervals are 95\% block-bootstrap confidence intervals (CIs; 10,000 resamples of 0.5° scoring blocks, conditional on the label draws), weighted by cell and by block as indicated in the key; the grey band marks ±0.5~cm.
In panel d, both spreading strategies lowered error in Canada by 2.5--3.2~cm at 50 and 100 labels and by 1.3--1.4~cm at 200 labels.
In the Lena Delta, block stratification changed error by +1.5 to +2.0~cm at 20--500 labels, with higher error under both weightings at 100 labels (+1.7~cm; depends on the split-independence assumption).
In Alaska the weightings disagreed at 50--200 labels (block-equal −1.1 to −1.7~cm, cell-weighted +0.28 to +0.74~cm).
Variance-based active selection had 0.69 to 2.8~cm higher error at 50--200 labels in Alaska and the Lena Delta (one of six contrasts depends on the split-independence assumption) and smaller gains than spreading in Canada.
In panel e, covariate spread lowered error in Canada at forty labels by 2.4~cm (95\% CI 0.5 to 3.4~cm; depends on the split-independence assumption); at ten labels no region had higher error under both weightings, but the cell-weighted interval of block stratification in the Lena Delta lay above zero (+1.9~cm).
The placement tests were designed after earlier results had been viewed and registered before running.
Polar stereographic projection centred on Canada, scale bar in panel a; coastlines, Natural Earth 50~m; Cartopy [미확인: 판].

\medskip
% ---- Figure 6 -------------------------------------------------------------
% Sources: C1 README 1.3, 2.1 (learner rows; five learners with agreeing CIs),
% 2.5 (post hoc counts 18 and 2 of 30; 17 + 10 + 3 targets), 2.7 (year-matched
% baseline); abstract bundle row for the main CatBoost learner (2.2459, rounded
% 2.2). Interval resamples: c 10,000 (v2 figure meta), d 1000 with calibration
% fixed (interval plan; one-stage columns chosen in figure spec 8.3), e 10,000.
% Placeholder: spec section 11.
\noindent Figure 6. Without target labels, 2 of 30 targets lost over 2~cm with physics pseudo-labels and 18 with direct ML.
Intervals in panel c are 95\% block-bootstrap confidence intervals (CIs; 10,000 resamples of 0.5° scoring blocks), weighted by cell and by block as indicated in the key, with the grey band at ±0.5~cm; bars in panels d and e are cell-weighted 95\% CIs (d, 1000 resamples, calibration fixed; e, 10,000 resamples).
(a, b) Error change relative to the source-coefficient Stefan model for direct ML (a) and direct ML trained with physics pseudo-labels (b) at each held-out target whose source excludes the parent region (17 targets); shared colour scale.
The title counts are a post hoc description of 30 targets (adding 10 sub-regions held out within their region and 3 expanded-label editions; threshold on point estimates).
(c) Error change of direct ML by learner and of two physics baselines, four-region mean.
The random forest and four neural networks configured by leave-one-source-region-out cross-validation had 2.4--5.2~cm higher error.
The other five learners, including the main CatBoost learner (2.2~cm; 95\% CI 1.1 to 3.4~cm; Holm-adjusted $P = 0.023$), gave the same verdict under the primary intervals, but this verdict depends on the split-independence assumption, as does the 1.0~cm lower error of Year-matched Stefan.
[XG: Fig 6c | 결과 열람 뒤 설계, 등록 이탈(WRAPUP 10) | 2.7 XG-1 행]
(d) Coverage of 90\% prediction intervals against mean width without labels for five interval scalings, four-region mean and regional values; the line marks nominal 0.90.
(e) Coverage of anchor plus residual ML intervals with 40 and 160 labels against width relative to zero labels; the grey band marks 0.85 to 0.95.
Polar stereographic projection (45°~W; scale true at 70°~N) with a Lambert azimuthal equal-area Tibetan inset; coastlines, Natural Earth 50~m; permafrost zones [미확인: 자료 이름과 판]; Cartopy [미확인: 판].

\medskip
% ---- Figure 7 -------------------------------------------------------------
% Sources: C4 README 2.1 (Alaska rows; the 1000-label row agrees under common
% resampling, 500 and all do not), 2.2 (Lena 23-24 splits), 2.6 (floors, 19%);
% C8 README 2.1, 2.2 (about 99%: 1 - 0.0086), 2.4 (three-region within-cell row;
% common-resampling verdict undetermined in the source table); Canadian floor
% line = 17.90 - 28.55 and 17.90 - 28.66 cm (figure spec 9.6); panel b mask and
% 343 locations from figure spec 9.3. Placeholder: spec section 11.
\noindent Figure 7. Within Alaska, residual ML gains over recalibrated Stefan were below 1~cm and came from climate-grid-scale bias correction.
Intervals are 95\% block-bootstrap confidence intervals (CIs; 10,000 resamples of 0.5° scoring blocks), weighted by cell and by block as indicated in the key; the grey band marks ±0.5~cm.
(a) Within-region tests (half of the 0.5° blocks scored, up to 25 splits): error change of Anchor + residual ML (cross-validated residual weight) and Direct ML relative to the Stefan model recalibrated on the same labels; Error floor is the covariate-conditional floor minus the recalibrated Stefan RMSE, and the Canadian line (−10.8 to −10.6~cm) is off the axis.
In Alaska with 1000 labels, Anchor + residual ML had 0.54~cm lower error (95\% CI 0.32 to 0.69~cm), removing 19\% of the squared error above the floor; the reductions with 500 and all labels (0.39 and 0.52~cm) did not hold under common resampling, and no difference was established in the Lena Delta or Canada.
(b) Alaska, all labels: error change per 0.5° scoring block, with label locations; hatching marks blocks with fewer than 10 scoring cells summed over the 25 splits.
(c) Error change split between and within ERA5-Land climate cells; in Alaska about 99\% of the squared-error reduction lay between cells (within cells −0.01~cm, equivalent), and in the three-region mean the within-cell predictions increased error by 0.09~cm (depends on the split-independence assumption).
(d) Recommended method per label interval at its error change relative to the source-coefficient Stefan model (Lena Delta and Canada transfer mean; 320 and 1000 labels untested), with placement and diagnosis steps below; selection from 40 labels rests on non-inferiority to a fixed recipe (Fig.~4c).
(e) [XC: Fig 7e | 결과 열람 뒤 설계, 재사용 지역의 재검정(비맹검 부분 포함) | 2.3 XC-1, XC-2 포레스트].
Panels a to d were designed after viewing the transfer results.
Polar stereographic projection centred on Alaska (scale true at 70°~N); coastlines, Natural Earth 50~m; Cartopy [미확인: 판].

% ---- Table 1 caption ------------------------------------------------------
% Sources: figure spec 10.2 to 10.4; D README 2 (row counts, source coefficient, Alaska share
% 78-94%); cell index 0.009 degrees from the v2 Table 1 caption. If a cell of the
% assembled table has no value, write "n.a." and add "n.a., not available." here
% (style guide 3.1). The table body is generated separately (Table1_data.tex).
\section*{Table caption}

\noindent Table 1. Target regions, labels and data sources.
Labels is the number of label rows; 1~km locations groups the rows by a 0.009° cell index; 0.5° blocks are the units of the label and scoring split; Splits is the number of valid block splits in the transfer design; $E_0$ is the Stefan coefficient (cm~(°C~d)$^{-1/2}$) fitted on the source pool when the region is held out; Alaska share of source cells is the fraction of source cells located in Alaska.
Main regions were used in earlier experiments and are retested here; Alaska is a reference target that is not averaged with them; new regions were added after the label tables had been fixed (licence-verified edition).
Sub-regions, targets with point estimates only and the North Atlantic group are listed in Supplementary Table~S1.
Access conditions are given under Data availability.
